\chapter*{\chapterstyle{I --- Introduction}}
\addcontentsline{toc}{section}{Introduction}

On appelle \textbf{structure algébrique} un ensemble muni d'une (ou plusieurs) opérations appelées \textbf{lois}, c'est l'étude de telles structures mathématiques, des relations entre celles-ci (que nous appeleront morphismes), et de leurs propriétés que nous appeleront \textbf{algèbre générale}.\<

Soit \(M\) un ensemble non-vide, on appelle \textbf{loi de composition interne} une opération binaire sur les éléments de \(M\) (qu'on notera temporairement \(\star\)) telle que:
\customBox{width=5cm}{
   \(\forall a, b \in M \; ; \; a \star b \in M\)
}
Soit \(E\) un ensemble non-vide, on appelle \textbf{loi de composition externe} une opération binaire entre un élément de \(E\) et un élément de \(M\) (qu'on notera temporairement \(\cdot\)) telle que:
\customBox{width=5cm}{
   \(\forall \lambda, a \in E \times M \; ; \; \lambda \cdot a \in M\)
}

On appelle \textbf{élément neutre} pour la loi un élément\footnote[1]{Dans la suite on notera l'élément neutre d'une structure \(M\) par \(e_M\)} \(e \in M\) tel que pour tout élément de \(a \in M\), on ait:
\[a \star e = e \star a = a\]

On appelle \textbf{inverse} pour la loi (si elle admet un élément neutre) un élément \(a^{-1} \in M\) tel que pour tout élément de \(a \in M\), on ait 
\[a \star a^{-1} = a^{-1} \star a = e\]

\subsection*{\subsecstyle{Monoides {:}}}
Soit \(M\) un ensemble qu'on munit d'une \textbf{loi de composition interne}, alors le couple \((M, \star)\) est appelé un \textbf{magma}, c'est la structure algébrique primitive la plus faible, en effet la seule contrainte étant que la loi soit interne.\<

On peut alors enrichir la structure de magma par les deux contraintes supplémentaires suivantes:
\begin{align*}
   &\bullet \;\; \text{La loi est \textbf{associative}.}\\
   &\bullet \;\; \text{Il existe \textbf{un élément neutre} pour la loi.}
\end{align*}
Cette structure plus riche, qu'on appelle \textbf{monoide} nous permet alors d'identifier des exemples remarquables:
\begin{align*}
   &\bullet \;\; \text{Les entiers naturels munis de l'addition forment un monoide.}\\
   &\bullet \;\; \text{L'ensemble des chaines de caractères muni de la concaténation forme un monoide.}
\end{align*}
Les éléments neutres respectifs de ces exemples sont \(0_\N\) et la chaine de caractère vide.

\subsection*{\subsecstyle{Morphismes {:}}}
Aprés avoir défini une structure, on peut alors définir les transformations qui \textbf{préservent cette structure}.\<

Soit \((M, \star)\) et \((N, \cdot)\) deux ensembles munis de la meme structure\footnote[2]{Si les structures présentent plusieurs lois, alors les morphismes doivent vérifier la compatibilité pour \textbf{toutes les lois}.}, \(x, y \in M\) et \(\varphi: M \rightarrow N\), alors \(\varphi\) est appelée \textbf{morphisme}, si elle vérifie:
\customBox{width=5cm}{
   \(
      \varphi(x \star y) = \varphi(x) \cdot \varphi(y)
   \)
}
Dans le cas particulier de structures qui requièrent l'existence d'un élément neutre, l'image de l'élement neutre de \(M\) par un morphisme doit etre l'élément neutre de \(N\).\<

En termes de vocabulaire, on définit alors: 
\begin{align*}
   &\bullet \;\; \text{\textbf{Les endormorphismes} commme les morphismes de \(M\) dans lui-meme.}\\
   &\bullet \;\; \text{\textbf{Les isomorphismes} commme les morphismes bijectifs.}\\
   &\bullet \;\; \text{\textbf{Les automorphismes} commme les morphismes bijectifs de \(M\) dans lui-meme.}
\end{align*}
Moralement, on comprends que:
\begin{center}
   \textit{Les morphismes préservent dans une certaine mesure la structure opératoire.}
\end{center}
En particulier:
\begin{center}
   \textit{Si il existe un isomorphisme entre deux structures, cela signifie que ces structures se comportent de la meme manière par rapport à leurs lois respectives.}
\end{center}
De manière plus subtile, si il existe un morphisme injectif d'une structure dans une autre, alors cela signifie qu'une partie de la seconde se comporte de la meme manière que la première. \<

Enfin, si il existe un morphisme surjectif d'une structure dans une autre, cela signifie qu'on peut regrouper des élements de la première structure et ces groupes d'éléments se comporteront commes les élements de la seconde\footnote[1]{Ces considérations plus avancées induiront l'idée générale des grands théorèmes de factorisation dans les groupes et les principaux morphismes canoniques.}.
\subsection*{\subsecstyle{Propriétés des morphismes {:}}}
Pour une structure donnée, si \(\varphi\) est un morphisme de \(E\) vers \(F\) et \(\psi\) est un morphisme de \(F\) vers \(G\), alors \(\psi \circ \varphi\) est un morphisme de \(E\) vers \(G\).\<

Si la structure admet un élément neutre pour la loi, alors:
\customBox{width=4cm}{
   \(\varphi(e_E) = e_F\)
}
Si la structure admet un symétrique pour la loi, alors:
\customBox{width=4cm}{
   \(\quad\;\varphi(x^{-1}) = \varphi(x)^{-1}\)
}
\subsection*{\subsecstyle{Exemples {:}}}
On peut considérer quelques exemples remarquables:
\begin{align*}
   &\bullet \;\; \text{\(\phi: n \in \N \longmapsto 2n \in 2\N\) est un morphisme (de monoide) de \((\N, +)\) dans \((2\N, +)\).}\\
   &\bullet \;\; \text{\(\phi: x \in \R \longmapsto \exp(x) \in \R^{+*}\) est un isomorphisme (de groupe) de \((\R, +)\) dans \((\R^{+*}, \times)\).}\\  
   &\bullet \;\; \text{\(\phi: M \in GL_{n}(\K) \longmapsto \det(M) \in \K^*\) est un morphisme (de groupe) de \((GL_{n}(\K), \times)\) dans \((\K, \times)\).}\\
   &\bullet \;\; \text{\(\phi: z \in \C \longmapsto \overline{z} \in \C\) est un automorphisme (de corps) de \((\C, +, \times)\) dans \((\C, +, \times)\).}
\end{align*}
\pagebreak
\subsection*{\subsecstyle{Sous-structures {:}}}

Une fois une structure algébrique définie sur \(E\), on peut alors s'intéresser aux parties de \(E\) qui conservent cette structure, on les appellera alors \textbf{sous-structures} de \(E\).\<

En particulier, on dira que \(F\) est une \textbf{sous-structure} de \(E\) (et on notera \(F < E\)) si elle vérifie:
\begin{itemize}
   \item La partie \(F\) est stable par les lois.
   \item Les éléments neutres\footnote[1]{\label{struct}\textbf{Si la structure impose leur existence}} appartient à \(F\) 
   \item Les inverses\footnoteref{struct} des éléments de \(F\) appartient à \(F\)
\end{itemize}

On peut alors montrer les propositions suivantes:
\begin{itemize}
   \item L'image et la préimage d'une sous-structure par un morphisme est une sous-structure.
   \item L'intersection d'une famille de sous-structures est une sous-structure
\end{itemize}
En particulier les noyaux de morphismes sont des sous-structures de la structure de départ.

\subsection*{\subsecstyle{Structures Quotients {:}}}
Dans le domaine ensembliste, on sait créer des ensembles quotients via une relation d'équivalence, on cherche par la suite à créer des ensembles quotients qui \textbf{conservent les propriétés de structure}, en particulier, on cherche des conditions sur une relation d'équivalence \(\sim\) pour que \(G/_\sim\) soit un groupe, un anneau, ou un corps.

On peut alors montrer qu'il faut et il suffit que \(\sim\) soit \textbf{compatible} avec les lois, ie que pour toute loi \(\star\), on ait:
\customBox{width=9cm}{
   \(
      x_1 \sim x_2 \text{ et } y_1 \sim y_2 \implies x_1 \star y_1 \sim x_2 \star y_2  
   \)
}

On peut alors montrer que \(G/_\sim\) conserve les propriétés de structure de \(G\).\<

On peut alors définir la \textbf{surjection canonique} qui est le morphisme \(\pi : E \rightarrow E /\sim\) qui a chaque élément associe sa classe d'équivalence pour la relation\footnote[1]{Non spécifique aux structures, c'est une application générale liée aux ensembles quotients.}. 

\chapter*{\chapterstyle{I --- Groupes}}
\addcontentsline{toc}{section}{Groupes}

Soit \(G\) un ensemble \textbf{non-vide} muni d'une loi de composition interne associative\footnote[1]{Dans la suite, la loi de composition des groupes sera notée multiplicativement sauf exceptions.} telle que:
\begin{align*}
   &\bullet \;\; \text{Il existe \textbf{un élément neutre} pour la loi.} \\
   &\bullet \;\; \text{Tout élément de \(G\) admet \textbf{un inverse} pour la loi.}
\end{align*}
Alors le couple \((G, \star)\) est appellé \textbf{groupe}. De plus si le groupe est \textbf{commutatif}, on dira alors que c'est un groupe \textbf{abélien}.\<

On appellera \textbf{ordre du groupe} le cardinal (potentiellement infini) de l'ensemble sous-jacent, noté \(|G|\).

\subsection*{\subsecstyle{Exemples {:}}}
On peut alors considérer plusieurs groupes remarquables:
\begin{align*}
   &\bullet \;\; \text{Les \textbf{entiers relatifs} muni de l'addition usuelle.}\\
   &\bullet \;\; \text{Les \textbf{isométries du plan} muni de la composition, on l'apelle le \textbf{groupe dihédral}.}\\   
   &\bullet \;\; \text{Les \textbf{matrices inversibles} muni de la multiplication, on l'apelle le \textbf{groupe linéaire}.}\\
   &\bullet \;\; \text{Les \textbf{bijections} sur un ensemble muni de la composition, on l'apelle le \textbf{groupe symétrique}.}
\end{align*}
\subsection*{\subsecstyle{Morphismes de groupes {:}}}
Soit \(G, H\) deux groupes et \(\varphi: G \rightarrow H\), l'existence d'un élément neutre nous permet de définir alors \textbf{le noyau d'un morphisme} par:
\[
   \text{Ker}(\varphi) := \Bigl\{ x \in G \; ; \; \varphi(x) = e_H \Bigl\}
\]
En particulier, on peut alors montrer:
\customBox{width=15cm}{
   Un morphisme est injectif si et seulement si son noyau est réduit à l'élément neutre.
}
\subsection*{\subsecstyle{Sous-groupes {:}}}
Les sous-structures dans le cas des groupes sont naturellement les \textbf{sous-groupes}. Un cas remarquable est celui du \textbf{sous-groupe engendré} par \(H\) qu'on note:
\customBox{width=10cm}{
   \(\langle H \rangle := \Bigl\{ h_1^{k_1}h_2^{k_2} \ldots h_n^{k_n}\; ; \; n \in \N \; , \; h_i \in H \; , \; k_i \in \Z \Bigl\}\)
}
\begin{center}
   \textit{C'est l'ensemble des combinaisons d'éléments de \(H\) et de leurs inverses.}
\end{center}
Une propriété fondamentale est que \(\langle H \rangle\) est un \textbf{opérateur de cloture} par la loi du groupe, ie c'est une application \textbf{idempotente, croissante et extensive}.\<

On peut alors considérer le sous-groupe engendré par un élément \(h \in H\), en effet on a:
\[
   \langle h \rangle := \bigl\{ h^{k} \; ; \; k \in \Z \bigl\}
\]
On peut alors définir \textbf{l'ordre d'un élément} comme étant l'ordre du sous-groupe engendré associé (potentiellement infini).\<

Ce sous-groupe permet de définir des groupes remarquables, en effet si un groupe est engendré par un unique élément, il est appelé \textbf{groupe cyclique} dont nous parleront plus loin dans ce chapitre.
\pagebreak
\subsection*{\subsecstyle{Action de groupe {:}}}
Soit \(E\) un ensemble et \(G\) un groupe noté multiplicativement, on peut définir une \textbf{action de groupe} de \(G\) sur \(E\) par une application de la forme:
\[
   \begin{aligned}
      G \times E &\longrightarrow E\\
      (g, x) &\longmapsto g \cdot x
   \end{aligned}
\]
Elle doit alors vérifier les propriétés suivantes:
\customBox{width = 10cm}{
   \begin{align*}
      &\textbf{Le neutre n'agit pas: }\forall x \in E \; ; \; e \cdot x = x\\
      &\textbf{Associativité mixte: }\forall g, g' \in G \; ; \; g' \cdot (g \cdot x) = (g'g) \cdot x
   \end{align*}
}

\underline{Exemple:} Considèrons le groupe \(\G_3\), il agit sur l'ensemble \(E = \left\{ A, B, C\right\}\) par l'action naturelle qui "mélange" l'ensemble \(E\), par exemple si on prends \(\sigma = (1, 2, 3)\), on associe \(\sigma(x)\) à chaque \(x \in E\).\+
Composer des permutations et faire agir la composée reviendra alors à faire agir les composantes successivement, et la permutation identité n'agit pas. 
\subsection*{\subsecstyle{Orbite {:}}}
On considère un ensemble \(E\) et un groupe \(G\) qui agit sur \(E\), alors on définit \textbf{l'orbite} d'un élément \(x \in E\) par:
\[
   O_x := \Bigl\{ y \in E \; ; \; \exists g \in G \, , \, g \cdot x = y\Bigl\}   
\]
\begin{center}
   \textit{L'orbite d'un élément est simplement l'ensemble des éléments atteint par l'action du groupe sur l'élément.}
\end{center}
\subsection*{\subsecstyle{Stabilisateur {:}}}
On peut aussi définir \textbf{le stabilisateur} d'un élément \(x \in E\) par:
\[
   St_x := \Bigl\{ g \in G \; ; \; g \cdot x = x\Bigl\}   
\]
\begin{center}
   \textit{Le stabilisateur d'un élément est simplement l'ensemble des éléments du groupe qui le laissent invariant.}
\end{center}
\subsection*{\subsecstyle{Classes {:}}}
Soit \(H < G\) et \(g \in G\), on appelle \textbf{classe à gauche}\footnote[1]{On définit de manière analogue les \textbf{classes à droite}.} de \(g\) suivant \(H\) l'ensemble:
\customBox{width=4.5cm}{
   \(gH := \Bigl\{ gh \; ; \; h \in H \Bigl\}\)
}
Ce sont les \textbf{classes d'équivalences} naturelles pour la relation \(a \sim b \Longleftrightarrow aH = bH\).\<

\underline{Exemple:} Considèrons le groupe \((\Z, \times)\), alors \(2\Z\) est la classe à gauche de \(2\) suivant \(\Z\). Les classes à gauche de la forme \(n\Z\) seront utiles pour définir les groupes quotients associés à \textbf{l'arithmétique modulaire}.
\subsection*{\subsecstyle{Sous-groupes normaux{:}}}
Soit \(H < G\) et \(g \in G\), on dit que \(H\) est un \textbf{sous-groupe normal} et on note \(H \lhd G\), si il vérifie:
\customBox{width=5cm}{
   \(
      \forall g \in G \; ; \; gH = Hg     
   \)
}
\begin{center}
   \textit{Un sous-groupe est normal si et seulement si toutes classe à gauche est égale à sa classe à droite.}
\end{center}
On peut alors montrer la propriété suivante:
\begin{center}
   \textbf{La préimage d'un sous-groupe normal par un morphisme est un sous-groupe normal\footnote[2]{Donc en particulier, les noyaux de morphismes sont toujours des sous-groupes normaux.}.}
\end{center}
\subsection*{\subsecstyle{Conjugaison{:}}}
On peut définir une action appellée \textbf{conjugaison} de \(G\) sur lui-meme par:
\[
   \begin{aligned}
      G \times G &\longrightarrow G\\
      (g, h) &\longmapsto ghg^{-1}
   \end{aligned}
\]
On peut alors montrer la caracterisation suivante pour \(H < G\):
\customBox{width=8cm}{
   \(
      H \unlhd G \Longleftrightarrow \forall g, h \in G \times H \; ; \; ghg^{-1} \in H   
   \)
}
\begin{center}
   \textit{Un sous-groupe est normal si et seulement si il est stable par conjugaison.}
\end{center}
\subsection*{\subsecstyle{Relations d'équivalence compatibles{:}}}
L'intéret fondamental des sous-groupes normaux et qu'on peut alors caractériser les relations d'équivalences \textbf{compatibles} avec la loi de groupe.\<

En effet si \(H \lhd G\), alors on définit:
\[
   x \sim y \Longleftrightarrow xH = yH
\]
On peut alors vérifier que la relation est compatible avec la loi du groupe et donc l'ensemble quotient suivant est bien un groupe pour la loi induite:
\[
   G/_\sim \underset{not.}{=} G/H
\]
\subsection*{\subsecstyle{Théorème de Lagrange{:}}}
Pour tout \textbf{groupe fini} et muni des outils précédemment définis (notamment des concepts de classes\footnote[1]{On peut partitionner \(G\) en classes suivant \(H\) qui ont toutes le meme cardinal que \(H\) et donc on a:\[
   |G| = |H| \times |G/H|\]}), on peut alors montrer la théorème fondamental suivant:
\begin{center}
   \customBox{width=11cm}{
      \textbf{L'ordre d'un sous-groupre divise l'ordre du groupe.}
   }
\end{center}

\chapter*{\chapterstyle{I --- Théorèmes d'isomorphismes}}
\addcontentsline{toc}{section}{Théorèmes d'isomorphismes}
On considère ici deux groupes \(G, F\), et \(H\) un sous-groupe normal de \(G\), on sait que \(H\) induit une relation d'équivalence compatible et donc qu'on peut définir le groupe quotient \(G/H\).\<

Soit \(\phi : G \longrightarrow F\) un morphisme, on peut alors montrer qu'il existe un unique morphisme \(\widetilde{\phi} : G/H \longrightarrow F\) tel que le diagramme soit commutatif\footnote[1]{Un \textbf{diagramme commutatif} est une collection d'objets et de morphismes tels tout les chemins (de composition) partant d'un objet vers un autre donnent le meme résultat (ie sont le meme morphisme).}:
\begin{center}
   \adjustbox{scale=1.5,center}{
      \begin{tikzcd}[column sep=large, row sep=large]
         G \arrow[d, "\pi", swap, two heads] \arrow[r, "\phi"] & F\\
         G/H \arrow[ru, "\widetilde{\phi}", swap, dashed]
      \end{tikzcd}
   }
\end{center}

\subsection*{\subsecstyle{Premier théorème d'isomorphisme {:}}}

\chapter*{\chapterstyle{I --- Groupes Symétriques}}
\addcontentsline{toc}{section}{Groupe Symétrique}
On appelle \textbf{groupe symétrique} et on note \(\mathfrak{S}_n\) le groupe des \textbf{permutations} de l'ensemble \(\inticc{1}{n}\) muni de la composition des applications.\<

On remarque alors aisément que l'ordre de \(\mathfrak{S}_n\) est \(n!\).\<

Soit \(\sigma \in \mathfrak{S}_n\) une permutation de \(\inticc{1}{n}\), alors c'est une fonction bijective sur cet ensemble. En particulier, sachant que l'ensemble est fini, c'est une fonction définie par cas qu'on note alors par commodité horizontalement dans un tableau:

\[
   \sigma =  \begin{pmatrix}
      1 & 2 & \ldots & n\\
      \sigma(1) & \sigma(2) & \ldots & \sigma(n)
   \end{pmatrix}
\]

\subsection*{\subsecstyle{Support {:}}}
On appelle alors \textbf{support} d'une permutation le complémentaire des points fixes de \(\sigma\), ie on a:
\[
   \text{Supp}(\sigma) := \bigl\{ i \in \N \; ; \; \sigma(i) \neq i \bigl\}   
\]
Une des propriétés fondamentale qu'on peut déduire de cette définition est que:
\customBox{width=10cm}{
   Deux permutations à supports disjoints \textbf{commutent}.
}
\subsection*{\subsecstyle{Cycles {:}}}
On appelle \textbf{k-cycle} une permutation \(\sigma\) telle qu'il existe \(k \geq 2\) et \(k\) éléments deux à deux distincts \(a_1, \ldots, a_n\) tels que:
\[
   \begin{cases}
      \forall i \in \inticc{1}{k-1} \; ; \; \sigma(a_i) = \sigma(a_{i+1}) \\
      \forall i \notin \inticc{1}{k-1} \; ; \; \sigma(a_i) = \sigma(a_{i})
   \end{cases} 
\]
\begin{center}
   \textit{Un k-cycle laisse fixe tout les éléments sauf pour une certaine famille \((a_i)\) pour laquelle chaque élément est envoyé sur le suivant.}
\end{center}
On peut alors noter un tel cycle par la notation suivante qui décrit tout les éléments affectés par la permutation:
\[
   \sigma = (a_1, \ldots, a_n)  
\]

Le cas particulier des \(2\)-cycles est intéressant, en effet un \(2\)-cycle \textbf{échange deux valeurs} de \(\inticc{1}{n}\), ils sont d'une importance particulière et on les appelle \textbf{transpositions}.\<

\underline{Exemple:} La permutation suivante est un 3-cycle:
\[
   \sigma =  \begin{pmatrix}
      1 & 2 & 3 & 4\\
      2 & 3 & 1 & 4
   \end{pmatrix} = (1 \;\, 2 \;\, 3)
\]
Si \(\sigma\) est un \(k\)-cycle et que \(a\) n'est pas un point fixe, alors on en déduit\footnote[1]{En effet par exemple \(\sigma(a)\) est la valeur suivante dans le cycle, et le cycle parcourt tout les points non-fixes par construction} que le support de \(\sigma\) est donné par:
\[
   \text{Supp}(\sigma) = \bigl\{a, \sigma(a), \sigma^2(a), \ldots, \sigma^{k-1}(a) \bigl\}
\]

\subsection*{\subsecstyle{Ordre {:}}}
On peut alors démontrer une propriété fondamentale de l'ordre des cycles:
\customBox{width=7cm}{
   \textbf{Un \(k\)-cycle est d'ordre \(k\)}.
}
En effet si on considère le sous-groupe engendré par un tel cycle, on remarque que pour tout élément \(a \in \inticc{1}{n}\) \(\sigma^{k}(a) = a\), donc \(\sigma^{k} = \text{Id}\).
\subsection*{\subsecstyle{Théorèmes de décomposition {:}}}
Une des problématiques principales à propos des groupes symétriques est la question de la \textbf{décomposition d'une permutation} en cycles. On peut en effet montrer que \textbf{toute permutation se décompose en produit de cycles à support disjoints}.\<

Par la suite, on peut alors constater directement que pour tout \(k\)-cycle \(\sigma = (a_1 \;\, \ldots \;\, a_k)\), on a:
\customBox{width=6cm}{
   \(\sigma = (a_1 \;\, a_2)(a_2 \;\, a_3)\ldots(a_{k-1} \;\, a_k)\)
}
Enfin, on conclura de ces deux propositions que \textbf{toute permutation se décompose en produit de transpositions}, ou en d'autres termes si on note \(\mathfrak{T}_n\) l'ensemble des transpositions:
\customBox{width=3cm}{
   \(\langle \mathfrak{T}_n \rangle = \mathfrak{S}_n\)
}
\subsection*{\subsecstyle{Algorithme de décomposition {:}}}
Dans cette section, on va présenter un algorithme permettant de décomposer une permutation quelconque en produit de transpositions.\+
On considère la permutation suivante:
\[
   \sigma = \begin{pmatrix}
      1 & 2 & 3 & 4\\
      4 & 3 & 1 & 2
   \end{pmatrix}  
\]
On va chercher à créer un point fixe pour la dernière valeur du tableau, on remarque alors qu'il suffit de composer \(\sigma\) par \((2 \;\, 4)\) et on obtient:
\[
   (2 \;\, 4)\sigma = \begin{pmatrix}
      1 & 2 & 3 & 4\\
      2 & 3 & 1 & 4
   \end{pmatrix}  
\]
On va chercher à créer un point fixe pour la troisème valeur du tableau, on remarque alors qu'il suffit de composer \((2 \;\, 4)\sigma\) par \((1 \;\, 3)\) et on obtient:
\[
   (1 \;\, 3)(2 \;\, 4)\sigma = \begin{pmatrix}
      1 & 2 & 3 & 4\\
      2 & 1 & 3 & 4
   \end{pmatrix}  
\]
Enfin il suffit de composer par \((1 \;\, 2)\) pour obtenir:
\customBox{width=5cm}{
   \((1 \;\, 2)(1 \;\, 3)(2 \;\, 4)\sigma = \text{Id}\)
}
Et donc il suffit alors de passer les transpositions à l'inverse pour obtenir que:
\[
   \sigma = (2 \;\, 4)(1 \;\, 3)(1 \;\, 2)
\]
\subsection*{\subsecstyle{Signature {:}}}
On appelle \textbf{inversion} par \(\sigma\) une paire \(\{i, j\} \in \inticc{1}{n}\) telle que:
\[
   i < j \text{ et } \sigma(i) > \sigma(j)  
\] 
On note alors \(\text{Inv}(\sigma)\) l'ensemble des inversions par \(\sigma\). On peut alors définir la \textbf{signature} de la permutation \(\sigma\) par:
\customBox{width=4cm}{
   \(\epsilon(\sigma) = (-1)^{|\text{Inv}(\sigma)|}\) 
}
Moralement:
\begin{center}
   \textit{La signature correspond à la parité du nombre d'inversions.}
\end{center}

\chapter*{\chapterstyle{I --- Groupes Cycliques}}
\addcontentsline{toc}{section}{Groupes Cycliques}
On appelle \textbf{groupe cyclique} un groupe \(G\) engendré par un unique élément qu'on notera \(g\). Le but de ce chapitre est de classifier ces groupes et d'identifier leurs caractéristiques.\<

Dans toute la suite, on utilisera le morphisme surjectif\footnote[1]{Car \(\langle g \rangle = G\) donc tout les éléments de \(G\) s'écrivent comme une puissance de \(g\).} suivant:
\[
   \begin{aligned}
      \phi: \Z &\longrightarrow G\\
      n &\longmapsto g^n
   \end{aligned}
\]

\subsection*{\subsecstyle{Cas infini {:}}}
Dans le cas ou \(G\) est infini, le morphisme \(\phi\) est injectif et donc on a la caractérisation suivante:
\customBox{width=2cm}{
   \(
      G \cong \Z
   \)
}
\begin{center}
   \textit{Il n'y a donc qu'un seul groupe cyclique d'ordre infini, celui des entiers naturels.}
\end{center}
\subsection*{\subsecstyle{Cas fini {:}}}
Dans le cas ou \(G\) est d'ordre \(n \in \N\), on peut utiliser le \textbf{premier théorème d'isomorphisme} pour montrer qu'il existe un isomorphisme:
\[
   \psi : \Z/\Ker{\phi} \longrightarrow G
\]
Et donc en particulier sachant que \(\Ker{\phi} = n\Z\), on a:
\customBox{width=3cm}{
   \(
      G \cong \Z/n\Z
   \)
}
\begin{center}
   \textit{Il n'y a donc qu'un seul groupe cyclique d'ordre n, celui des classes de congruences modulo n.}
\end{center}
\subsection*{\subsecstyle{Indicatrice d'Euler {:}}}
On définit \textbf{la fonction indicatrice d'Euler} par:
\[
   \begin{aligned}
      \varphi: \N &\longrightarrow \N\\
      n &\longmapsto n \prod_{p|n}{\left(1-\frac{1}{p}\right)}
   \end{aligned}
\]
Le produit se faisant sur tout les diviseurs premiers distincts de \(n\). L'utilité de cette fonction vient de la propriété suivante:
\customBox{width=13cm}{
   \textbf{Le nombre d'entiers inférieurs à \(n\) et premiers avec \(n\) est de \(\varphi(n)\).}
}

\underline{Exemple:} \(\varphi(30) = \varphi(2\times3\times5) = 30\times\left(1-\frac{1}{2}\right)\left(1-\frac{1}{3}\right)\left(1-\frac{1}{5}\right) = 30\times\frac{1}{2}\times\frac{2}{3}\times\frac{4}{5} = 8\)
\subsection*{\subsecstyle{Propriétés {:}}}
On peut alors de démontrer les propriétés suivantes:

\begin{center}
   \begin{itemize}
      \item Tout groupe cyclique est \textbf{abélien}.
      \item Tout \textbf{sous-groupe} d'un groupe cyclique est cyclique.
   \end{itemize}
\end{center}

\chapter*{\chapterstyle{I --- Anneaux}}
\addcontentsline{toc}{section}{Anneaux}
Soit \(A\) un ensemble \textbf{non-vide} muni de deux lois de composition internes associatives notées \(+, \times\) telles que:
\begin{align*}
   &\bullet \;\; (A, +) \text{ soit un groupe commutatif.} \\
   &\bullet \;\; \text{La loi \(\times\) est associative.}\\
   &\bullet \;\; \text{La loi \(\times\) est distributive sur la loi \(+\).}\\
   &\bullet \;\; \text{Il existe \textbf{un élément neutre} pour la loi \(\times\).}
\end{align*}
Alors le triplet \((A, +, \times)\) est appellé \textbf{anneau}. Si la loi multiplicative est \textbf{commutative}, on dira alors que c'est un anneau commutatif.
\subsection*{\subsecstyle{Exemples {:}}}
On peut alors considérer plusieurs anneaux remarquables:
\begin{itemize}
   \item Les \textbf{entiers relatifs} muni des opérations usuelles.
   \item Les \textbf{fonctions continues} muni de la somme et du produit.
   \item Les \textbf{polynomes}\footnote{\label{coeff}A coefficients dans un anneau} muni de la somme et du produit.
   \item Les \textbf{matrices}\footnoteref{coeff} muni de la somme et du produit.
\end{itemize}

\subsection*{\subsecstyle{Propriétés Algébriques{:}}}
Pour deux éléments \(a, b \in A\) qui commutent, on a \textbf{la formule du binome de Newton}:
\[
   (a + b)^n = \sum_{k=0}^{n}\binom{n}{k} a^k b^{n-k}   
\]

\subsection*{\subsecstyle{Sous-anneaux {:}}}
Les sous-structures dans le cas des groupes sont naturellement les \textbf{sous-anneaux}. Un cas remarquable est celui du \textbf{sous-anneau engendré} par \(H\) qu'on note:
\customBox{width=10cm}{
   \(\langle H \rangle := \Bigl\{ \sum_{k=1}^{n} \pm h_1^{k_1}h_2^{k_2} \ldots h_n^{k_n}\; ; \; n \in \N \; , \; h_i \in H \; , \; k_i \in \N \Bigl\}\)
}
\begin{center}
   \textit{C'est l'ensemble des produits et sommes d'éléments de \(H\) et de leurs inverses pour la loi de groupe.}
\end{center}
Une propriété fondamentale est que \(\langle H \rangle\) est un \textbf{opérateur de cloture} par la loi du groupe, ie c'est une application \textbf{idempotente, croissante et extensive}.
\subsection*{\subsecstyle{Idéaux{:}}}
On appelle \textbf{idéal} tout sous groupe additif de \(A\) qui soit stable par multiplication (à droite et à gauche) par n'importe quel élément de l'anneau.
\begin{center}
   \textit{Les idéaux jouent alors le meme role que les sous-groupes normaux, ie on peut quotienter par ceux-ci.}
\end{center}
On peut alors montrer les propriétés suivantes:
\begin{itemize}
   \item La préimage d'un idéal par un morphisme est un idéal\footnote[2]{Donc en particulier, les noyaux de morphismes sont toujours des idéaux.}.
   \item L'intesection d'une famille d'idéaux est un idéal.
\end{itemize}

\subsection*{\subsecstyle{Inversibles {:}}}
On dit qu'un élément \(x \in A^*\) est \textbf{inversible} à droite\footnote[2]{On définit de meme les inversibles à gauche} si et seulement si il existe \(y \in A\) tel que:
\[
   xy = 1
\]
Si un élément est inversible bilatère, on dira alors simplement qu'il est inversible. L'ensemble des inversibles d'un anneau forme un groupe pour la loi multiplicative qu'on note \(\mathbb{U}(A)\).
\subsection*{\subsecstyle{Diviseurs de zéro {:}}}
On dit qu'un élément \(x \in A^*\) est \textbf{un diviseur de zéro} à droite\footnote[3]{On définit de meme les diviseurs de zéro à gauche} si et seulement si il existe \(y \in A^*\) tel que:
\[
   yx = 0
\]
\underline{Exemple}: Dans l'anneau \(\mathcal{M}_2(\R)\), la matrice \(\begin{pmatrix}1 & 1 \\ 0 & 0\end{pmatrix}\) est un diviseur de zéro car \(\begin{pmatrix}1 & 1 \\ 0 & 0\end{pmatrix}\begin{pmatrix}1 & 0 \\ -1 & 0\end{pmatrix} = 0\)
\subsection*{\subsecstyle{Nilpotents {:}}}
On dit qu'un élément \(a \in A\) est \textbf{nilpotent} si et seulement si il existe \(n \in \N\) tel que:
\[
   a^n = 0
\]
En particulier les nilpotents sont donc des diviseurs de zéro. \<

\underline{Exemple}: Dans l'anneau \(\mathcal{M}_2(\R)\), la matrice \(A = \begin{pmatrix}0 & 1 \\ 0 & 0\end{pmatrix}\) est nilpotente car \(A^2 = 0\)
\subsection*{\subsecstyle{Caractéristique {:}}}
On définit la caractéristique d'un anneau non-nul par:
\[
   \text{char}(A) := \min \left\{ n \in \N \; ; \;\underbrace{1 + \ldots + 1}_{\text{\(n\) sommandes}} = 0 \right\}
\]
Une autre formulation serait simplement que:
\begin{center}
   \textit{La caractéristique d'un anneau est l'ordre (additif) de l'unité multiplicative.}
\end{center}
\subsection*{\subsecstyle{Anneaux intègres {:}}}
On appelle \textbf{anneau intègre} tout anneau \(A\) non-nul \textbf{commutatif} et \textbf{sans diviseurs de zéro}. En particulier, dans un anneau intègre, on a alors la propriété qu'un produit est nul si et seulement si \textbf{l'un des facteurs est nul}.
\subsection*{\subsecstyle{Anneaux à PGCD {:}}}
\subsection*{\subsecstyle{Anneaux Factoriels {:}}}
\subsection*{\subsecstyle{Anneaux Principaux {:}}}
\subsection*{\subsecstyle{Anneaux Euclidiens {:}}}
On appelle \textbf{anneau Euclidiens} tout anneau \(A\) principal qui possède une \textbf{division euclidienne}. Dans un tel anneau, on peut alors faire \textbf{de l'arithmétique} comme dans l'anneau des entiers naturels.
\subsection*{\subsecstyle{Schéma heuristique des structures d'anneaux {:}}}
Pour mieux visualiser la hierarchie des différents types d'anneaux, on peut représenter la structure logique sous la forme de la suite d'implications suivantes:
\begin{center}
   \customBox{width = 16cm}{
      \textbf{Euclidien} \(\implies\) \textbf{Principal} \(\implies\) \textbf{Factoriel} \(\implies\) \textbf{PGCD} \(\implies\) \textbf{Intégre} \(\implies\) \textbf{Commutatif}
   }
\end{center}

\chapter*{\chapterstyle{I --- Corps}}
\addcontentsline{toc}{section}{Corps}
Soit \(A\) un anneau dont tout les éléments sauf \(0\) sont inversibles. Alors on dit que \(A\) est \textbf{un corps}.

\subsection*{\subsecstyle{Exemples {:}}}
On peut alors considérer plusieurs corps remarquables:
\begin{itemize}
   \item Les \textbf{réels} muni des opérations usuelles.
   \item Les \textbf{quaternions}\footnote[1]{C'est un exemple de corps non commutatif} muni des opérations usuelles.
   \item Les \textbf{corps finis} \(\mathbb{F}_p = \Z/p\Z\) pour \(p\) premier.
   \item Les \textbf{nombres constructibles} à la règle et au compas.
\end{itemize}

\chapter*{\chapterstyle{I --- Arithmétique dans \(\Z\)}}
\addcontentsline{toc}{section}{Arithmétique dans \(\Z\)}

\subsection*{\subsecstyle{Division Euclidienne {:}}}
Soit \(a, b \in \Z \times \Z^*\), on montre qu'il existe un unique couple \((q, r) \in \Z \times \N\) avec \(r < |b|\) tel que:

\customBox{width=4cm}{
   \(
      a = bq + r
   \)
}
On appelle alors cette décomposition \textbf{la division euclidienne} de \(a\) par \(b\).

\subsection*{\subsecstyle{Plus grand diviseur commun {:}}}

Soit \(a, b \in \Z\) non simultanément nuls, alors le pgcd est l'entier \(d\) qui vérifie:
\customBox{width=5cm}{
   \begin{align*}
      &\bullet \;\; d \, | \, a \text{ et } d \, | \, b \\
      &\bullet \;\; d' \, | \, a \text{ et } d' \, | \, b \implies d' \leq d
   \end{align*}
}
Alors on l'appelle \textbf{plus grand diviseur commun} de \(a\) et de \(b\) et on le note \(a \wedge b\). Pour le trouver en pratique, on peut utiliser l'algorithme d'Euclide. \<

Une caractérisation utile permet de montrer que si \(d\) est un diviseur commun à \(a\) et \(b\), alors \(d = \text{pgcd}(a , b)\) si et seulement si pour tout diviseur commun \(d'\) de \(a\) et \(b\):
\customBox{width=2cm}{
   \(
      d' \, | \, d   
   \)
}

\subsection*{\subsecstyle{Théorème de Bézout {:}}}
Soit \(a, b\in \Z^2\), on peut montrer qu'il existe deux entiers \(n_1, n_2 \in \Z^2\) tels que:
\customBox{width=5cm}{
   \(
      n_1a+n_2b = a \wedge b
   \)
}
\begin{center}
   \textit{
       Il existe donc une combinaison linéaire (à coefficients entiers) de \(a, b\) qui donne leur PGCD.
   }
\end{center}

\subsection*{\subsecstyle{Lemme de Gauss {:}}}
Soit 3 entiers \(a, b, c \in \Z\), alors\footnote[1]{Partir d'une combinaison linéaire donnée par Bézout.}:
\customBox{width=5.5cm}{
   \[
      \begin{cases} 
         a \, | \, bc \\
         a \wedge b = 1
      \end{cases} \implies a \, | \, c
   \]
}